\scanpage{7}
\uline{Démonstration du Lemme 1} Il sera une conséquence du\\
\uline{Lemme 2} \inserted{(Mumford)} Soit $B$ schéma abélien sur un corps [alg. clos.] $k$,\\
et $L$ un Module inv. sur $B$ alg. équivalent à zéro, \deleted{\unsure{Alors}}\\
et non isomorphe à $\uline{\mathcal O}_B$. Alors on a
\[
H^i(X,L)=0\quad \deleted{\boxed{\unsure{\text{si}}\ L\not\simeq\mathcal O_B}}\quad\text{pour }\deleted{\text{tout}}\ i
\]
\deleted{et $H^0(X,L)=\boxed{k\ \text{si }L\simeq\mathcal O_B}$}.

\deleted{\uline{Démonstration du Lemme 2} L'assertion sur les $H^0$ est triviale,\\
et il reste à prouver que $H^i(X,L)=0$ \unsure{si $i>0$}.}

Nous allons ici admettre le lemme 2, et nous \unsure{montrer|montrerons}\\
comment en \deleted{\unsure{conclure}} \inserted{déduire} le lemme 1. Nous voyons déjà\\
grâce à lemme 2 que $\mathbf Rp_*(W)$ est concentré sur la\\
section unité de $B^*$. Nous allons définir un\\
homomorphisme canonique (où on pose $e_{B^*}=e$)
\[
\mathbf Rp_*(W)\longrightarrow\deleted{\unsure{e_*}(\mathcal O_S)[-n]}\ e_*(\uline{\mathcal O}_S)\otimes_S\tau_{B^*}[-n]
\]
ou, ce qui revient au $=$, \inserted{(comme $n$ est la dim relative de $B/S$, donc $R^ip_*(W)=0$ si $i>n$)} un homomorphisme\\
canonique
\[
R^np_*(W)\longrightarrow\deleted{\unsure{e_*(\mathcal O_S)}}\ e_*(\tau_{B^*}\deleted{\unsure{[-n]}})
\]
ou encore un hom.
\[
\unsure{\mathrm{eq.}|e_{B^*}}\quad e^*(R^np_*(W))\longrightarrow\deleted{\mathcal O_S}\ \tau_{B^*}.
\]
\diagram{\begin{tikzcd}[ampersand replacement=\&,column sep=large,row sep=large]
B^*\times_S B \arrow[d,"p"'] \& B \arrow[l,"e'"'] \arrow[d,dash,"q"] \\
B^* \& S \arrow[l,"e"]
\end{tikzcd}}
Or en dim. $n$ on a \inserted{formule de} changement de base,\\
et \inserted{comme $e'^*(W)\simeq\uline{\mathcal O}_B$,} on est ramené à définir un hom.
\[
R^nq_*(\uline{\mathcal O}_B)\longrightarrow\deleted{\mathcal O_S}\ \tau_{B^*}.
\]
Or on sait que l'on a un isom canonique
\[
R^nq_*(\mathcal O_B)\simeq\overset{n}{\Lambda}R^1q_*(\mathcal O_B)\simeq\overset{n}{\Lambda}t_{B^*}
\]
où $t_{B^*}$ est le faisceau tangent de $B^*$ le long de $e$.\\
On \deleted{\unsure{fera correspondre}} \inserted{prend alors} \uline{l'hom identique} $\overset{n}{\Lambda}t_{B^*}\simeq\tau_{B^*}$.
\mspage{6}
